\vfill\pagebreak
\section{Sharing what can change: \texttt{alias}}
\label{sec:mutability:alias}

In Listing~\ref{lst:mutability:share}, \token{seen} shares the elements of
\token{primes} but cannot change them. A second variable that may change them
too is refused, if it is declared as \token{seen} was.

\begin{lstlisting}[style=coloredverbatimError, escapechar=@, caption=Two variables that may change the same elements]
use std::io;

fn main() {
  let dmut scores = copy [12, 15, 9];
  let dmut same = @\hb{scores}@; // error
  same[2] = 10;
  println(scores, " ", same);
}
\end{lstlisting}

The compiler says \errmsg{E4045}{discarding the constant property is
  prohibited}, and the note under it gives the reason: \textit{implicit alias of
  type mut [mut i32] is prohibited, it implicitly gives up on borrowings of
  mutable values}. An \textit{alias} is a second name for the same elements.
The declaration would create one silently, and the line \token{same[2] = 10;}
would then change \token{scores}, which it does not name.

A reader of a long function cannot keep in mind every pair of variables that
share elements. Ymir therefore asks for such a pair to be created in writing,
with the keyword \token{alias} in front of the value shared.

\lstinputlisting[style=coloredverbatim, caption={A second name for the same elements, \textit{examples/mutability/alias.yr}}, label=lst:mutability:alias]{examples/mutability/alias.yr}
\vspace{-10pt}%
\begin{lstlisting}[style=bashVerb]
[12, 15, 10] [12, 15, 10]
\end{lstlisting}

\token{alias scores}, on line~5, gives the length and the address of
\token{scores}, as a plain assignment would. The keyword changes nothing to what
is copied; it records that the program means the two variables to share
elements that both may change. A reader who looks for the places where
\token{scores} may change without being named finds them by looking for
\token{alias scores}.

\subsection{Only where it is needed}

\token{alias} is required where a value whose elements may change is given to a
variable whose elements may change too. Elsewhere, it is refused, so that it
keeps marking only these places.

\begin{lstlisting}[style=coloredverbatimError, escapechar=@, caption=An alias to a constant]
use std::io;

fn main() {
  let dmut scores = copy [12, 15, 9];
  let seen = @\hb{alias}@ scores; // error
  println(scores, " ", seen);
}
\end{lstlisting}

The compiler says \errmsg{E4247}{aliasing the value of type mut [mut i32] to
  create a constant borrowing is prohibited}. \token{seen} cannot change the
elements, and \token{let seen = scores;} is the declaration to write. The
keyword is refused as well in front of a variable whose elements are constant,
such as one declared \token{let mut}: it has no right to change them, and
cannot give that right to another variable.

\subsection{Tuples and arrays that hold slices}

A tuple is stored in its variable, like an array, and its assignment copies its
fields. A field that is a slice is copied as a slice is, length and address, and
its elements are then shared.

\begin{lstlisting}[style=coloredverbatim]
let dmut team = (1, copy [12, 15, 9]);
let dmut other = alias team;
other.0 = 2;
other.1[0] = 20;
println(team, " ", other);
\end{lstlisting}
\vspace{-10pt}%
\begin{lstlisting}[style=bashVerb]
(1, [20, 15, 9]) (2, [20, 15, 9])
\end{lstlisting}

\token{other.0 = 2;} changes the first field of \token{other}, which is a copy:
\token{team.0} is unchanged. \token{other.1[0] = 20;} writes
into the elements that both second fields designate. \token{alias} is needed
because of that second field, whose elements both variables may change. An
array of slices behaves the same way.
